\chapter{Chỉ đường lại}
\markboth{Chỉ đường lại}{Show the Way Again}


\section{Không biết bắt đầu từ đâu thì bắt đầu từ chỗ dễ nhất.}
\begin{truyen}{Không biết bắt đầu từ đâu thì bắt đầu từ chỗ dễ nhất.}{Start with the easiest step.}
\chuhoa{K}{hi nản, học sinh thường thấy mọi thứ rối thành một khối. Người thầy tốt không bắt các em nuốt cả ngọn núi, mà chỉ ra viên đá đầu tiên để bước qua.}
\textit{(When discouraged, students see everything as one giant mess. A good teacher does not ask them to swallow the mountain, but points to the first small stone to step over.)}
\end{truyen}

\begin{baihoc}
Đường ra khỏi bế tắc thường bắt đầu bằng việc nhỏ nhất, không phải lớn nhất.
\textit{(The way out of confusion usually begins with the smallest step, not the biggest.)}
\end{baihoc}

\begin{ghinhoanh}
\textbf{Short English to remember:}

Start small. Start now.

\end{ghinhoanh}

\begin{tuhocnhanh}
1. Việc nhỏ nhất em có thể làm ngay là gì? \textit{(What is the smallest thing you can do right now?)}

2. Em có đang đợi mình đủ hứng mới bắt đầu không? \textit{(Are you waiting to feel ready before starting?)}
\end{tuhocnhanh}
\ngancach


\section{Đừng nghĩ về cả tháng — nghĩ về buổi tối nay trước đã.}
\begin{truyen}{Đừng nghĩ về cả tháng — nghĩ về buổi tối nay trước đã.}{Shrink the timeline.}
\chuhoa{N}{ghĩ quá xa làm học sinh mệt trước khi bắt đầu. Khi thời gian được thu nhỏ lại thành một buổi tối, lòng người bớt hoảng và công việc trở nên khả thi hơn.}
\textit{(Thinking too far ahead exhausts students before they start. When time is reduced to just one evening, the heart calms down and the task becomes possible.)}
\end{truyen}

\begin{baihoc}
Chia nhỏ thời gian là cách giảm sợ hãi và tăng hành động.
\textit{(Breaking time into smaller parts reduces fear and increases action.)}
\end{baihoc}

\begin{ghinhoanh}
\textbf{Short English to remember:}

Do tonight first.

\end{ghinhoanh}

\begin{tuhocnhanh}
1. Nếu chỉ lo cho tối nay, em sẽ làm gì? \textit{(If you only think about tonight, what would you do?)}

2. Việc gì em đang làm quá lớn trong đầu mình? \textit{(What have you made too big in your mind?)}
\end{tuhocnhanh}
\ngancach


\section{Thầy không chỉ em đường tắt — thầy chỉ em đường đúng.}
\begin{truyen}{Thầy không chỉ em đường tắt — thầy chỉ em đường đúng.}{The right road matters.}
\chuhoa{C}{ó những lúc học sinh nản đến mức chỉ muốn thoát nhanh. Nhưng đường tắt sai có thể làm các em lạc xa hơn. Điều các em cần là con đường đúng và đủ rõ để đi tiếp.}
\textit{(Sometimes students are so discouraged they only want a quick escape. But the wrong shortcut can lead them farther astray. What they need is the right path, made clear enough to follow.)}
\end{truyen}

\begin{baihoc}
Đi đúng chậm còn hơn đi sai thật nhanh.
\textit{(Going slowly in the right direction is better than rushing the wrong way.)}
\end{baihoc}

\begin{ghinhoanh}
\textbf{Short English to remember:}

Right first. Fast later.

\end{ghinhoanh}

\begin{tuhocnhanh}
1. Em có đang muốn tìm đường tắt không? \textit{(Are you looking for shortcuts right now?)}

2. Đường đúng của em lúc này là gì? \textit{(What is the right path for you now?)}
\end{tuhocnhanh}
\ngancach


\section{Khi rối, em đừng làm tất cả — hãy làm điều quan trọng nhất trước.}
\begin{truyen}{Khi rối, em đừng làm tất cả — hãy làm điều quan trọng nhất...}{Prioritize when confused.}
\chuhoa{R}{ối trí thường đến khi học sinh muốn ôm hết mọi thứ. Một người thầy tốt sẽ chỉ ra: không phải việc nào cũng cần làm ngay, nhưng luôn có một việc quan trọng nhất.}
\textit{(Confusion often comes when students try to hold everything at once. A good teacher points out that not everything must be done now, but there is always one most important thing.)}
\end{truyen}

\begin{baihoc}
Ưu tiên đúng làm học sinh bớt loạn và bớt nản.
\textit{(Good prioritizing helps students feel less chaotic and less discouraged.)}
\end{baihoc}

\begin{ghinhoanh}
\textbf{Short English to remember:}

Do the important thing first.

\end{ghinhoanh}

\begin{tuhocnhanh}
1. Việc quan trọng nhất của em hôm nay là gì? \textit{(What is your most important task today?)}

2. Em đang phân tán vì những việc nào? \textit{(What things are scattering your attention?)}
\end{tuhocnhanh}
\ngancach


\section{Nếu em lạc, quay về điều căn bản nhất.}
\begin{truyen}{Nếu em lạc, quay về điều căn bản nhất.}{Back to basics.}
\chuhoa{C}{ó những lúc học sinh nản vì cố hiểu cái khó khi nền chưa chắc. Quay lại điều căn bản không phải lùi bước; đó là dựng lại nền nhà.}
\textit{(Sometimes students are discouraged because they try to master hard things without a solid base. Returning to fundamentals is not moving backward; it is rebuilding the foundation.)}
\end{truyen}

\begin{baihoc}
Căn bản vững thì đường học mới không làm em hoảng.
\textit{(A solid foundation keeps the learning path from feeling terrifying.)}
\end{baihoc}

\begin{ghinhoanh}
\textbf{Short English to remember:}

Back to basics is not failure.

\end{ghinhoanh}

\begin{tuhocnhanh}
1. Phần căn bản nào em đang hổng? \textit{(Which basic part are you weak in?)}

2. Em có dám quay lại học lại nó không? \textit{(Do you dare to go back and relearn it?)}
\end{tuhocnhanh}
\ngancach


\section{Thầy không cần em đi hết đường trong hôm nay.}
\begin{truyen}{Thầy không cần em đi hết đường trong hôm nay.}{One day at a time.}
\chuhoa{N}{hiều học sinh mệt vì tự đặt cho mình đòi hỏi quá lớn trong một ngày. Câu nói đúng sẽ kéo các em về nhịp thực tế hơn: hôm nay chỉ cần đi một đoạn.}
\textit{(Many students are exhausted because they demand too much from one day. The right sentence pulls them back to reality: today, just walk one stretch.)}
\end{truyen}

\begin{baihoc}
Đường dài cần nhịp bền, không cần nước rút mỗi ngày.
\textit{(A long road needs steady rhythm, not daily sprints.)}
\end{baihoc}

\begin{ghinhoanh}
\textbf{Short English to remember:}

Do not finish everything today.

\end{ghinhoanh}

\begin{tuhocnhanh}
1. Em đang đòi mình điều gì quá mức trong hôm nay? \textit{(What are you demanding too much of yourself today?)}

2. Một đoạn đường vừa sức của em hôm nay là gì? \textit{(What is a manageable stretch for today?)}
\end{tuhocnhanh}
\ngancach


\section{Lối ra không phải lúc nào cũng là cố hơn — đôi khi là đổi cách.}
\begin{truyen}{Lối ra không phải lúc nào cũng là cố hơn — đôi khi là đổi ...}{Change the method.}
\chuhoa{H}{ọc sinh thường nghĩ bế tắc nghĩa là mình phải cố gấp đôi. Nhưng đôi khi, điều cần không phải cố hơn mà là làm khác đi cho đúng hơn.}
\textit{(Students often think being stuck means they must try twice as hard. But sometimes the answer is not more force, but a better method.)}
\end{truyen}

\begin{baihoc}
Cố gắng mà sai cách chỉ làm người ta mệt thêm.
\textit{(Effort with the wrong method only creates more exhaustion.)}
\end{baihoc}

\begin{ghinhoanh}
\textbf{Short English to remember:}

Do not only try harder. Try wiser.

\end{ghinhoanh}

\begin{tuhocnhanh}
1. Em đang cố nhiều hay đang đổi cách đủ chưa? \textit{(Are you only trying harder, or changing your method too?)}

2. Có cách nào em nên thử khác đi? \textit{(What method should you try differently?)}
\end{tuhocnhanh}
\ngancach


\section{Mỗi lần em hỏi lại là mỗi lần em bớt lạc hơn.}
\begin{truyen}{Mỗi lần em hỏi lại là mỗi lần em bớt lạc hơn.}{Questions guide the way.}
\chuhoa{H}{ỏi lại không phải là phiền. Với học sinh đang nản, việc dám hỏi lại chính là dấu hiệu các em vẫn còn muốn đi tiếp.}
\textit{(Asking again is not a burden. For discouraged students, daring to ask again is a sign they still want to continue.)}
\end{truyen}

\begin{baihoc}
Biết hỏi lại là biết tự kéo mình ra khỏi mù mờ.
\textit{(Knowing how to ask again helps you pull yourself out of confusion.)}
\end{baihoc}

\begin{ghinhoanh}
\textbf{Short English to remember:}

Ask again. Get clearer.

\end{ghinhoanh}

\begin{tuhocnhanh}
1. Em có hay ngại hỏi lại không? \textit{(Do you often hesitate to ask again?)}

2. Điều gì làm em sợ khi phải hỏi? \textit{(What scares you about asking questions?)}
\end{tuhocnhanh}
\ngancach


\section{Cứ làm rõ bước kế tiếp, rồi con đường sẽ hiện ra dần.}
\begin{truyen}{Cứ làm rõ bước kế tiếp, rồi con đường sẽ hiện ra dần.}{Clarity grows step by step.}
\chuhoa{C}{on đường học không phải lúc nào cũng sáng từ đầu đến cuối. Nhưng nếu bước kế tiếp đủ rõ, học sinh sẽ có can đảm để nhích tiếp.}
\textit{(The learning road is not always clear from start to finish. But when the next step is clear enough, students regain courage to move forward.)}
\end{truyen}

\begin{baihoc}
Sự rõ ràng nhỏ tạo ra lòng can đảm lớn hơn ta tưởng.
\textit{(Small clarity creates more courage than we often realize.)}
\end{baihoc}

\begin{ghinhoanh}
\textbf{Short English to remember:}

See the next step. Then move.

\end{ghinhoanh}

\begin{tuhocnhanh}
1. Bước kế tiếp rõ nhất của em là gì? \textit{(What is your clearest next step?)}

2. Em cần ai giúp làm rõ nó? \textit{(Who can help you make it clearer?)}
\end{tuhocnhanh}
\ngancach


\section{Lạc một đoạn không có nghĩa là em đã mất cả con đường.}
\begin{truyen}{Lạc một đoạn không có nghĩa là em đã mất cả con đường.}{A detour is not the end.}
\chuhoa{C}{ó học sinh chỉ vì đi sai vài bước mà tưởng mình hỏng cả chặng đường. Nhưng đi lạc một đoạn không xóa hết những gì đã học, cũng không cắt đứt cơ hội quay lại.}
\textit{(Some students make a few wrong moves and think the whole journey is ruined. But a short detour does not erase what they have learned, nor does it remove the chance to return.)}
\end{truyen}

\begin{baihoc}
Sai hướng một lúc không đáng sợ bằng việc vì xấu hổ mà không dám quay lại.
\textit{(A temporary wrong direction is less dangerous than being too ashamed to turn back.)}
\end{baihoc}

\begin{ghinhoanh}
\textbf{Short English to remember:}

A detour is not the end.

\end{ghinhoanh}

\begin{tuhocnhanh}
1. Em đang lạc ở đoạn nào, và em có thể quay lại từ đâu? \textit{(Where are you off track, and from where can you return?)}

2. Điều gì khiến em ngại sửa hướng của mình? \textit{(What makes you hesitate to correct your direction?)}
\end{tuhocnhanh}
\ngancach